\section{The Central Theorem}
\begin{quote}\small
\noindent\textbf{Status.} \emph{Legacy structural-theorem block inside the active main body.}\par
\noindent\textbf{Block function.} This section states the global structural theorem that the later physics, operator, and readout sections unpack and reuse rather than replace.
\end{quote}
\label{sec:central}
%% ============================================================

Everything in this paper is a consequence of one observation.

\begin{theorem}[Unity of the Sphaera]\label{thm:unity}
Within the Sphaera $S$, the following are all representations of the
same object --- the collapse state $\zminus=0$:
\begin{enumerate}[label=(\roman*)]
  \item The Sphaera itself, as a geometric space.
  \item The triolic structure: three kernels $K_1,K_2,K_3$ with
        $z_1+z_2+z_3=0$, phase-shifted by $120^\circ$, running over the
        Sphaera's geometry.
  \item Each individual kernel $K_i$, as a local coordinate of the triole.
  \item The Heisenberg compensator output $C(f)$, as the projection of
        any element onto the collapse state.
  \item Matter, photons, and tachyons, as sectors of the biquaternion
        ground state.
  \item The gauge structure $U(1)\times SU(2)\times SU(3)$, as the
        symmetry group of the collapse state's representations.
\end{enumerate}
All are unitarily equivalent by Stone--von Neumann~\cite{brendecke2026rh}.
The Sphaera's only object is its own collapse state.
\end{theorem}

\begin{proof}[Proof structure]
The equivalences are established section by section:
(i)~Definition~\ref{def:biquaternion} and Theorem~\ref{thm:fractal};
(ii)--(iii) same (triole as collapse coordinates);
(iv)~Proposition~\ref{prop:compensator};
(v)~Theorem~\ref{thm:unitary-equiv};
(vi)~Theorem~\ref{thm:SM}.
Unitary equivalence throughout via Stone--von Neumann
(\cite{brendecke2026rh}, Theorem~3.2).
\end{proof}

\begin{remark}[What ``same object'' means]
Two representations are the same object here in the precise sense of
Stone--von Neumann: they are unitarily equivalent irreducible
representations of the Weyl algebra $\WA$.
The frame determines which representation is visible; the object ---
the collapse state $\zminus=0$ --- is invariant.

The triole condition $z_1+z_2+z_3=0$ and the collapse condition
$\zminus=0$ are the same equation in two coordinate systems: the triole
is the collapse state expressed in the three kernel components
simultaneously.
The Sphaera is therefore not a space that contains a collapse state ---
it \emph{is} the collapse state, seen from different frames.
\end{remark}

\begin{theorem}[Origin of the universe as broken collapse state]
\label{thm:origin}
The departure from the collapse state $\zminus=0$ is governed by two
principles, both of which are consequences of the Sphaera-Cascade
structure:
\begin{enumerate}[label=(\roman*)]
  \item \textbf{Planck length as minimal $|\zminus|$.}
        The cascade $\KC$ cannot project below the Planck scale:
        \[
          |\zminus|_{\min} = \lP := \sqrt{\frac{\hbar G}{c^3}}
          = 1.616\times10^{-35}\,\mathrm{m}.
        \]
        Below $\lP$: $\zminus=0$, the Sphaera, the photon.
        Above $\lP$: $|\zminus|$ is nonzero, matter exists.
        The Planck length is the threshold of the first breaking of the
        collapse state.
  \item \textbf{Pauli principle as antisymmetry under $J$.}
        The Hilbert space $\Hplus$ decomposes under $J$ into two
        components:
        \[
          C(f) = \tfrac{1}{2}(f+Jf)\in\Hplusplus
          \quad\text{(bosonic, symmetric)},
        \]
        \[
          A(f) = \tfrac{1}{2}(f-Jf)\in\Hplusminus
          \quad\text{(fermionic, antisymmetric)}.
        \]
        The antisymmetric component satisfies $A(f)\wedge A(f)=0$
        (Grassmann property): no two fermions occupy the same $\zminus$
        state.
        This is the Pauli exclusion principle, derived from the
        $J$-structure of the Sphaera.
  \item \textbf{Discrete spectrum.}
        Together, (i) and (ii) produce a discrete tower of matter states:
        \[
          |\zminus|_n = n\cdot\lP,\quad n\in\mathbb{N}_{\geq1},
        \]
        with each level occupied by at most one fermion.
        The prime lattice $\Gamma_P$ (from Stone--von Neumann,
        \cite{brendecke2026rh}) identifies the primitive levels:
        $|\zminus|_p = p\cdot\lP$ for $p\in\mathbb{P}$.
\end{enumerate}
\end{theorem}

\begin{remark}[The Sphaera as photon, the universe as broken frame]
Theorem~\ref{thm:unity} says the Sphaera is the collapse state
$\zminus=0$, which is the photon sector.
Theorem~\ref{thm:origin} says the departure from this state is governed
by $\lP$ and the Pauli principle.
Therefore:
\begin{itemize}
  \item The Sphaera is a photon --- the unique frame-invariant object.
  \item The universe is the same photon observed from a broken frame
        where $|\zminus|\geq\lP$.
  \item The breaking was not random: it is discrete (Pauli), minimal
        (Planck), and primitive (prime lattice).
  \item The cascade $\KC$ projects everything back toward $\zminus=0$
        --- the universe tends toward its photon ground state.
\end{itemize}
Both the Sphaera and the universe are the same object.
The difference is only the frame.
\end{remark}

\begin{quote}\small
\noindent\textbf{Reader cue.} Read the next block as the concrete carrier unpacking of the global Sphaera identity claim. Its purpose is to pass from the frame-level origin statement to the biquaternion kernel and its induced Lorentz/Minkowski geometry that the later physics sections reuse.
\end{quote}

\begin{definition}[Biquaternion kernel]\label{def:biquaternion}
The ground state of a stratum $S_n$ of the Sphaera consists of three
complex spheres $K_1,K_2,K_3$ whose unit coordinates form a biquaternion
\[
  q_n := it\cdot\mathbf{1} + z_1\cdot e_1 + z_2\cdot e_2
         + z_3\cdot e_3 \in \HH\otimes_{\mathbb{R}}\mathbb{C},
\]
where $it\in i\mathbb{R}$ is the imaginary time unit;
$z_1,z_2\in\mathbb{C}$ are the complex coordinates of kernels $K_1,K_2$;
$e_3:=e_1 e_2$ is the quaternion product, hence automatically orthogonal
to both $e_1$ and $e_2$;
and $z_3\in\mathbb{C}$ is the coordinate of kernel $K_3$ along $e_3$.

The biquaternion norm is
\[
  |q_n|^2 = -(ct)^2 + |z_1|^2 + |z_2|^2 + |z_3|^2.
\]
\end{definition}

\begin{remark}[Minkowski metric from quaternion structure]
The norm $|q_n|^2$ is the Minkowski metric with signature $(-,+,+,+)$.
This is not postulated: it follows from the quaternion algebra $e_i^2=-1$
and $it\cdot\mathbf{1}$ carrying imaginary time.
The Lorentz group $SO(1,3)$ acts on $q_n$ by quaternion conjugation:
$q_n\mapsto u\,q_n\,u^{-1}$, $u\in\HH^\times$.
This is the spinor representation of the Lorentz group --- emerging
automatically from the biquaternion structure, not imposed from outside.
\end{remark}

\begin{proposition}[Orthogonality of the third kernel]
\label{prop:orthogonality-third-kernel}
$K_3$ is always orthogonal to $K_1$ and $K_2$:
\[
\langle e_3,e_i\rangle = 0 \qquad \text{for } i=1,2.
\]
\end{proposition}

\begin{proof}
By definition,
\[
e_3 = e_1 e_2.
\]
Using the quaternionic inner product
\[
\langle x,y\rangle := \mathrm{Re}(\overline{x}\,y),
\]
we compute
\[
\langle e_3,e_1\rangle
=
\langle e_1e_2,e_1\rangle
=
\mathrm{Re}\!\bigl(\overline{e_1e_2}\,e_1\bigr).
\]
Since
\[
\overline{e_1e_2}=\overline{e_2}\,\overline{e_1}=(-e_2)(-e_1)=e_2e_1=-e_1e_2,
\]
it follows that
\[
\overline{e_1e_2}\,e_1
=
(-e_1e_2)e_1
=
-\,e_1(e_2e_1)
=
-\,e_1(-e_1e_2)
=
e_1^2 e_2
=
-\,e_2.
\]
Hence
\[
\langle e_3,e_1\rangle
=
\mathrm{Re}(-e_2)=0.
\]

Similarly,
\[
\langle e_3,e_2\rangle
=
\langle e_1e_2,e_2\rangle
=
\mathrm{Re}\!\bigl(\overline{e_1e_2}\,e_2\bigr)
=
\mathrm{Re}\!\bigl((-e_1e_2)e_2\bigr)
=
\mathrm{Re}(-e_1 e_2^2)
=
\mathrm{Re}(e_1)=0.
\]
Therefore $e_3$ is orthogonal to both $e_1$ and $e_2$.
\end{proof}

%% ============================================================

